\documentclass[10pt,aspectratio=169]{beamer}
\usepackage[portuguese]{babel}
\usepackage{fontspec,unicode-math}
\setmainfont{Noto Sans}
\setsansfont{Noto Sans}
\setmathfont{STIX Math}
\usepackage{amsmath}
\usepackage{xcolor}
\definecolor{InsperRed}{HTML}{a01924}
\definecolor{Ink}{HTML}{232529}
\setbeamercolor{frametitle}{fg=Ink}
\setbeamercolor{structure}{fg=InsperRed}
\setbeamercolor{normal text}{fg=Ink,bg=white}
\setbeamerfont{frametitle}{size=\large,series=\bfseries}
\setbeamerfont{title}{size=\LARGE,series=\bfseries}
\setbeamertemplate{navigation symbols}{}
\setbeamertemplate{itemize item}{\color{InsperRed}$\bullet$}
\setbeamertemplate{footline}{%
  \hbox{\begin{beamercolorbox}[wd=.84\paperwidth,ht=2.7ex,dp=1.2ex,leftskip=.7cm]{normal text}\scriptsize Inteligência Artificial para Decisão · Insper · 2027-1\end{beamercolorbox}%
  \begin{beamercolorbox}[wd=.16\paperwidth,ht=2.7ex,dp=1.2ex,rightskip=.7cm plus1fil]{normal text}\hfill\scriptsize\insertframenumber\, /\, \inserttotalframenumber\end{beamercolorbox}}}
\setbeamersize{text margin left=12mm,text margin right=12mm}
\begin{document}
\begin{frame}[t]{}
\vspace{14mm}\textcolor{InsperRed}{\bfseries AULA 02}\par\vspace{4mm}{\LARGE\bfseries Embeddings e similaridade semântica}\par\vspace{6mm}{\large Marketing Analítico}\par\vspace{12mm}{\color{gray}Paulo C. Marques F. · 2027-1}
\end{frame}

\begin{frame}[t]{Objetivos e verificação}
\small
\begin{itemize}
\item Interpretar embeddings, medidas de similaridade e agrupamentos exploratórios, produzindo uma análise replicável de avaliações e marcas.
\item Resolver manualmente os exemplos do capítulo, identificando hipóteses e consequências.
\item Especificar a aplicação no Codex gráfico e confrontar ao menos um resultado com evidências.
\end{itemize}
\end{frame}

\begin{frame}[t]{Organização do encontro}
\small
\begin{itemize}
\item Apresentação dos fundamentos e resolução de problemas quantitativos no papel.
\item Laboratório supervisionado e testes produzidos por prompts na interface gráfica do Codex, sem terminal.
\item Atividade individual com dados diferentes e depósito no GitHub antes do quiz seguinte.
\end{itemize}
\end{frame}

\begin{frame}[t]{Representação vetorial de textos}
\small
\begin{itemize}
\item Um modelo de embeddings associa a cada texto um vetor .
\item A representação é aprendida e depende do modelo, do idioma e do pré-processamento.
\item Não se deve assumir que cada coordenada represente diretamente um atributo de mercado.
\item Vetores são instrumentos para comparar descrições, mas a relação entre semântica e geometria não é perfeita: sinônimos, ironia, negação e mudança de domínio podem alterar as distâncias.
\end{itemize}
\end{frame}

\begin{frame}[t]{Representação vetorial de textos — continuação}
\small
\begin{itemize}
\item No caso de avaliações, é necessário definir a unidade estatística.
\item Um comentário, uma avaliação e um cliente não são a mesma observação.
\item Um consumidor prolífico pode produzir dezenas de textos e dominar a distribuição dos comentários; isso não implica que domine a distribuição dos consumidores.
\item Para segmentação de clientes, precisamos de identificadores, critérios de agregação e amostragem apropriada.
\end{itemize}
\end{frame}

\begin{frame}[t]{Produto interno, norma e similaridade}
\small
\begin{itemize}
\item Para vetores não nulos , definimos A similaridade cosseno pertence ao intervalo para vetores reais não nulos.
\item Ela é invariável a multiplicação por constante positiva: .
\item Já a distância euclidiana entre e é .
\item Portanto, alinhamento não significa proximidade absoluta nem intensidade idêntica. Exemplo resolvido. Para , e , temos , e .
\end{itemize}
\end{frame}

\begin{frame}[t]{Produto interno, norma e similaridade — continuação}
\small
\begin{itemize}
\item A diferença entre a ordenação por cosseno e por distância deve ser explicitada antes de interpretar um gráfico.
\end{itemize}
\end{frame}

\begin{frame}[t]{Produto interno, norma e similaridade — formulação matemática}
\small
\vspace{3mm}\begin{center}\small\[\langle u,v\rangle=\sum_{j=1}^d u_jv_j,\qquad
\|u\|_2=\sqrt{\sum_{j=1}^d u_j^2},\qquad
\cos(u,v)=\frac{\langle u,v\rangle}{\|u\|_2\|v\|_2}.\]\end{center}
\vspace{3mm}\small A similaridade cosseno pertence ao intervalo para vetores reais não nulos.
\end{frame}

\begin{frame}[t]{Agregação, agrupamentos e projeção}
\small
\begin{itemize}
\item O centroide de vetores é .
\item Ele minimiza a soma das distâncias euclidianas quadráticas dentro de um grupo, mas pode ser distorcido por observações extremas.
\item Um algoritmo de agrupamento é exploratório: a existência de dois clusters numéricos não prova a existência de dois segmentos comerciais.
\item É indispensável examinar estabilidade, conteúdo dos grupos e composição dos dados.
\end{itemize}
\end{frame}

\begin{frame}[t]{Agregação, agrupamentos e projeção — continuação}
\small
\begin{itemize}
\item Projetar embeddings para duas dimensões pode facilitar a inspeção, mas sempre requer atenção à distorção.
\item Se a distância original satisfaz e a figura apresenta o contrário, o mapa não preservou nem a ordenação local desse trio.
\item Para posicionamento de marcas, atributos, preços e participação de mercado devem provir de fontes próprias; não são inferíveis diretamente de uma nuvem de comentários.
\end{itemize}
\end{frame}

\begin{frame}[t]{Aplicação à decisão}
\small
\begin{itemize}
\item Na análise de marcas, a pergunta relevante pode ser: ``Quais atributos estão semanticamente associados às avaliações de cada marca, e com que incerteza?''.
\item A aplicação pode recuperar avaliações exemplares, calcular centroides e localizar observações discrepantes.
\item Deve mostrar também quantas avaliações e quantos clientes geraram cada resultado, a distribuição de notas e casos discordantes.
\item Um gráfico sem essas informações pode ter grande apelo visual e pequeno valor analítico.
\end{itemize}
\end{frame}

\begin{frame}[t]{Exercício para resolver no papel 1}
\small
\small\textbf{2.1 O mapa de três marcas}\par\vspace{4mm}Os protótipos de percepção A, B e C são vetores , e . Calcule , e . Por que a marca mais próxima pela direção não precisa ser a mais próxima pela distância euclidiana?
\end{frame}

\begin{frame}[t]{Exercício para resolver no papel 2}
\small
\small\textbf{2.2 Uma campanha que não muda a direção}\par\vspace{4mm}Um comentário tem embedding . Outro tem . Calcule o cosseno e as normas. Um sistema conclui que são duas opiniões com intensidade idêntica porque o cosseno é máximo. Qual é a falha dessa conclusão?
\end{frame}

\begin{frame}[t]{Laboratório e caso aplicado}
\small
\begin{itemize}
\item Área de aplicação: Marketing Analítico.
\item Comparar manualmente dois vetores bidimensionais e formular uma hipótese sobre uma frase próxima e outra distante.
\item Pedir ao Codex embeddings locais, matriz de cossenos, recuperação dos três vizinhos e visualização exploratória com rótulo de marca.
\item Testar cossenos para frases idênticas, matriz simétrica e casos de negação; comparar duas escolhas simples de agrupamento.
\end{itemize}
\end{frame}

\begin{frame}[t]{Prompt de desenvolvimento — exemplo}
\small
\small\textbf{Instrução fornecida ao Codex:}\par\vspace{5mm}Construa um explorador local de embeddings para o conjunto de 48 avaliações. Mostre cossenos, vizinhos, grupos e textos originais; registre o modelo e parâmetros. Evite API convencional e não interprete grupos de comentários como perfis demográficos.
\end{frame}

\begin{frame}[t]{Critérios de aceitação e validação}
\small
\begin{itemize}
\item Cossenos simétricos e próximos de 1 na diagonal de vetores não nulos; modelo identificado; grupos acompanhados de exemplos; nenhuma afirmação sobre market share ou número de clientes sem denominador apropriado.
\item Inspecionar argumentos, saídas e rastros das ferramentas quando houver chamadas externas.
\item Comparar uma saída com a resolução manual do exemplo conceitual.
\end{itemize}
\end{frame}

\begin{frame}[t]{Atividade individual em casa}
\small
\begin{itemize}
\item Base distinta com 36 avaliações fictícias de três marcas de protetor solar; inclua comentários sobre preço, conforto e duração, alguns com negação e um duplicado.
\item Resolver as questões de interpretação, com cálculos e evidência documental.
\item Depositar no GitHub os arquivos solicitados; a entrega habilita o quiz, sem valer nota.
\end{itemize}
\end{frame}

\begin{frame}[t]{Preparação para a avaliação escrita}
\small
\begin{itemize}
\item Exercícios do caderno: 2.1 O mapa de três marcas; 2.2 Uma campanha que não muda a direção; 2.3 Centroide e voz minoritária.
\item Formalizar o raciocínio e interpretar resultados; não será exigida escrita manual de código.
\item No quiz seguinte serão avaliados conceitos desta aula e da atividade feita em casa.
\end{itemize}
\end{frame}

\end{document}
